\documentclass[10pt, a4paper]{article}
\usepackage{preamble}

\title{Statistical Inference II \\
    \large Prereading Notes}
\author{Luke Phillips}
\date{October 2025}

\begin{document}

\maketitle

\newpage

\tableofcontents

\newpage

\section{Multivariate probability}

\begin{definition}[Random vector]
    If $X_1, \dotsc, X_p$ is a collection of real-valued random variables,
    then
    \[
    \mbf{X} = \begin{bmatrix}
        X_1 \\ \vdots \\ X_p
    \end{bmatrix} = (X_1, \dotsc, X_p) ^ T,
    \]
    is a $p$-dimensional random vector in $\R ^ p$,
    and $\mbf{x} = (x_1, \dotsc, x_p) ^ T$ is a particular realisation of $\mbf{X}$.
\end{definition}

\begin{definition}[PDF of a random vector]
    The pdf of a random vector $\mbf{X} = (X_1, \dotsc, X_p) ^ T$ is the joint probability density function
    (also just density or pdf)
    $f(X_1, \dotsc, X_p)$ from $\R ^ p$ to $\R_{0}^{+} = [0, \infty)$ such that:
    \begin{enumerate}[label = \roman*)]
        \item $f(\mbf{x}) \geq 0$ for all $\mbf{x} \in \R ^ p$.

        \item $\int_{\R ^ p}f(\mbf{x})\,d\mbf{x} = \int_{-\infty}^{\infty}\dotsi\int_{-\infty}^{\infty}f(X_1, \dotsc, X_p)\,dX_1\dotsc dX_p = 1$.

        \item For any subset $\mathcal{A}$ of $\R ^ p = \P(\mbf{X} \in \mathcal{A}) = \int_{\mathcal{A}}f(\mbf{x})\,d\mbf{x}$.
    \end{enumerate}
\end{definition}

\begin{definition}[CDF of a random vector]
    Similarly,
    the joint cdf is defined as
    \[
    F(\mbf{x}) = \P(X_1 \leq x_1, \dotsc, X_p \leq x_p) = \int_{-\infty}^{x_1} \dotsi \int_{-\infty}^{x_p}f(x_1', \dotsc, x_p')\,dx_p'\dotsc dx_1',
    \]
    and we can recover the pdf by
    (partial)
    differentiation
    \[
    f(\mbf{x}) = \frac{\partial ^ p}{\partial x_1, \dotsc, \partial x_p}F(\mbf{x}).
    \]
\end{definition}

\subsection{Key operations and relations}
\begin{example}
    Consider the bivariate
    $(p = 2)$ pdf:
    \[
    f_{X, Y}(x, y) = \begin{cases}
        \lambda ^ 2e ^ {-\lambda y}, &0 < x < y,\text{ and } \lambda > 0 \\
        0, &\text{otherwise}.
    \end{cases}
    \]
    Is this a valid pdf?

    \begin{solution}
        We know $\lambda ^ 2 > 0$ and $e ^ x > 0$ for all $x \in \R$,
        so we have $f_{X, Y}(x, y) \geq 0$ for all $(x, y) \in \R ^ 2$ hence we have positivity.

        Now we need to check that the integral integrates to $1$ over $\R ^ 2$:
        \begin{align*}
            \int_{-\infty}^{\infty}\int_{-\infty}^{\infty}f_{X, Y}(x, y)\,dxdy &= \int_{0}^{\infty}\int_{0}^{y}\lambda ^ 2e ^ {-\lambda y}\,dxdy \\
            &= \int_{0}^{\infty}\lambda ^ 2e ^ {-\lambda y}\int_{0}^{y}\,dxdy \\
            &= \int_{0}^{\infty}\lambda ^ 2e ^ {-\lambda y}y\,dy \\
            &= \left[-(\lambda y + 1)e ^ {-\lambda y}\right]_{0}^{\infty} \\
            &= 0 - (-1) = 1.
        \end{align*}
        Hence,
        the integral equals $1$ and we have positivity so this is a valid pdf.
    \end{solution}
\end{example}

\begin{definition}[Marginal pdf]
    A sub-collection $(X_1, \dotsc, X_q)$ of $(X_1, \dotsc, X_p)$,
    where $q < p$,
    has a marginal pdf defined as
    \[
    f(x_1, \dotsc, x_q) = \int_{-\infty}^{\infty}\dotsi\int_{-\infty}^{\infty}f(x_1, \dotsc, x_p)\,dx_{q + 1}\dotsc dx_p. 
    \]
    For discrete random variables,
    the same definition holds with integrals replaced by sums.
\end{definition}

\begin{definition}[Conditional pdf]
    If we know the values $X_{q + 1} = x_{q + 1}, \dotsc, X_p = x_p$,
    but not the values of $X_1, \dotsc, X_q$,
    then the $X_1, \dotsc, X_q$ are still random variables and have a conditional pdf
    \[
    f(x_1, \dotsc, x_q\mid x_{q + 1}, \dotsc, x_p) = \frac{f(x_1, \dotsc, x_q, x_{q + 1}, \dotsc, x_p)}{f(x_{q + 1}, \dotsc, x_p)}.
    \]
\end{definition}

\begin{definition}[Independence]
    A collection of random variables $(X_1, \dotsc, X_p)$ is said to be
    (mutually)
    independent if and only if their joint pdf factorises as
    \[
    f(x_1, x_2, \dotsc, x_p) = f_1(x_1)f_2(x_2)\dotsc f_p(x_p),
    \]
    for all $(x_1, x_2, \dotsc, x_p) \in \R ^ o$.
    We will sometimes write this as $X_1 \independent X_2 \dotsc \independent X_p$ or $X_i \independent X_j\ \forall i \neq j$,
    where the symbol $\independent$ means 'is independent of'.

    An equivalent statement,
    is that the conditional densities are equal to the marginal densities:
    \[
    f(x_i \mid x_1, \dotsc, x_{i - 1}, x_{i + 1}, \dotsc, x_p) = f(x_i),
    \]
    for all $i = 1, \dotsc, p$.
\end{definition}

\begin{definition}[Random sample,
    IID]
    The random variables $X_1, \dotsc, X_p$ are called independent and identically distributed
    (iid)
    random variables with pdf or pmf $f(x)$ if they are all independent of each other,
    and the marginal pdf or pmf of each $X_i$ is the same function $f(x)$.
    This means that the joint pdf simplifies to
    \[
    f(x_1, x_2, \dotsc, x_p) = f(x_1)f(x_2) \dotsc f(x_p) = \prod_{i = 1}^{p}f(x_i).
    \]
    Alternatively,
    $X_1, \dotsc, X_p$ are called a random sample of size $n$ from the distribution
    (or population)
    $f(x)$.
\end{definition}

\newpage

\section{Expectations,
means,
and variances}

\subsection{Expectation of a scalar function of multiple random variables}

\begin{definition}[The honesty condition]
    The honesty condition for a pdf
    (or pmf)
    says that the integral
    (or sum)
    of a valid pdf over all possible values must equal $1$.
\end{definition}

\begin{remark}
    When working with probability densities if we can rearrange our integral
    (or part of it)
    to be the integral of a known valid pdf over its entire range,
    then we can immediately state that it must integrate to one.
\end{remark}

\subsection{Expectation and variance of a random vector}

\begin{definition}[Mean vector]
    The expectation $\E(\mbf{X})$ of a random vector $\mbf{X} = (X_1, \dotsc, X_p)$ is defined as:
    \[
    \E(\mbf{X}) = \int\mbf{x}f(\mbf{x})\,d\mbf{x} = \begin{bmatrix}
        \int_{-\infty}^{\infty}x_1f(\mbf{x})\,d\mbf{x} \\
        \vdots \\
        \int_{-\infty}^{\infty}x_pf(\mbf{x})\,d\mbf{x}
    \end{bmatrix} = \begin{bmatrix}
        \E(X_1) \\ \vdots \\ \E(X_p)
    \end{bmatrix} = \mbf{\mu},
    \]
    which is simply the $p$-dimensional vector whose components are the expectations of the individual components of $\mbf{x}$ and is known as the mean vector,
    $\mbf{\mu}$,
    of $\mbf{X}$.
\end{definition}

\begin{definition}[Variance matrix]
    The variance matrix of the random vector $\mbf{X}$ is defined to be the $p \times p$ matrix such that,
    for $i, j = 1, \dotsc, p$,
    the $(i, j)$ element of the matrix is $\Cov(X_i, X_j)$:
    \[
    \Var(\mbf{x}) = \E((\mbf{X} - \mbf{\mu})(\mbf{x} - \mbf{\mu}) ^ T) = \E(\mbf{X}\mbf{X} ^ T) - \mbf{\mu}\mbf{\mu} ^ T = \begin{bmatrix}
        \Sigma_{11} & \dotsi & \Sigma_{1p} \\
        \vdots & \ddots & \vdots \\
        \Sigma_{p1} & \dotsi & \Sigma_{pp}
    \end{bmatrix} = \mbf{\Sigma},
    \]
    where the components of the matrix comprise the individual variances and covariances of the random variables:
    \begin{align*}
        \Sigma_{jj} &= \E(X_j ^ 2) - \E(X_j) ^ 2 = \Var(X_j) = \sigma_j ^ 2 \\
        \Sigma_{ij} &= \E(X_iX_j) - \E(X_i)\E(X_j) = \Cov(X_i, X_j) = \Cov(X_j, X_i) = \Sigma_{ji}.
    \end{align*}
\end{definition}

\begin{definition}[Covariance and correlation]
    The covariance,
    $\Cov(X_i, X_j)$ between a pair of random variables $X_i$ and $X_j$ is a measure of the strength of linear association between the two variables,
    defined as:
    \begin{align*}
        \Cov(X_i, X_j) &= \E((X_i - \E(X_i))(X_j - \E(X_j))) = \Cov(X_j, X_i) \\
        &= \sigma_i\sigma_j\rho_{ij},
    \end{align*}
    where $\sigma_i, \sigma_j$ are the standard deviations of $X_i$ and $X_j$ respectively,
    and $\rho_{ij}$ is the correlation between $X_i$ and $X_j$.

    The correlation $\rho_{ij}$ between $X_i$ and $X_j$ is a unit-less measure of association which quantifies the extent to which two random variables are linearly related.
    The correlation is such that $\rho_{ij} \in [-1, 1]$,
    with values further from zero indicating stronger association.
    If $X_1$ and $X_2$ are independent then $\rho_{ij} = 0$,
    but the reverse is not true in general.
\end{definition}

\subsection{Data,
sample means and variances}

\begin{definition}[Data matrix]
    The measurement of the $j$th variable on the $i$th individual is denoted $x_{ij}$,
    and is stored as the $(i, j)$th element of the $n \times p$ data matrix
    \[
    \mbf{X} = \begin{bmatrix}
        x_{11} & x_{12} & \dotsi & x_{1p} \\
        x_{21} & x_{22} & \dotsi & x_{2p} \\
        \vdots & \vdots & \ddots & \vdots \\
        x_{n1} & x_{n2} & \dotsi & x_{np} \\
    \end{bmatrix}.
    \]
\end{definition}
The $i$th row of $\mbf{X}$ is the observation vector of all $p$ variables for individual $i$,
$\mbf{x}_i$,
and then the $j$th column of $\mbf{X}$ corresponds to all of the observation of variable $j$ across the $n$ individuals.

\end{document}